% \iffalse
%<*driver>
\ProvidesFile{seqenseigne-legacy.dtx}%
[2025/01/11 v1.0-dev seqenseigne – Module legacy (compatibilité ascendante)]
%</driver>
% \fi
%
% =============================================================
% seqenseigne-legacy.dtx
% Module : LEGACY
% Rôle   :
%   - macros et environnements historiques, conservés
%     pour que l'ancien code continue de compiler
%   - wrappers et alias pour compatibilité
%   - environnements avec anciens noms
%
%   → À terme ce module pourra être vidé puis retiré.
% =============================================================
%
%<*package>
\NeedsTeXFormat{LaTeX2e}%
\ProvidesPackage{seqenseigne-legacy}%
[2025/01/11 v1.0-dev seqenseigne legacy module]

%% ------------------------------------------------------------------
%% ---- Traitement des options --------------------------------------
%% ------------------------------------------------------------------
\newif\ifdocstd\docstdtrue
\DeclareOption{ltxdoc}{\docstdfalse}%
\ProcessOptions\relax

% longtblr (utilisé dans titreFiche) requiert la bibliothèque longtable de tabularray.
% tabularray lui-même est déjà chargé par seqenseigne-core.
% \UseTblrLibrary{longtable}%

%% ------------------------------------------------------------------
%% ---- Alias de variables ------------------------------------------
%% ------------------------------------------------------------------
% Note: \codeSequence est supprimé. Utiliser \seqCodeSequence ou \seqSetCodeSequence.

%% ------------------------------------------------------------------
%% ---- Alias d'environnements --------------------------------------
%% ------------------------------------------------------------------

% Les anciens alias d'environnements (colEnum, colItem, corrige, prerequisN2,
% exercice, serieExos) sont supprimés. Utiliser les noms courants :
%   seqColEnum, seqColItem, seqCorrige, seqExercice, seqSerieExos.

%% ------------------------------------------------------------------
%% ---- Macros obsolètes mais conservées ----------------------------
%% ------------------------------------------------------------------

% Ligne pour chaque élève (fiche de suivi)
% \tabucline remplacé par \hline (compatible longtblr)
\newcommand{\seq@ligne@eleve}[1]{%
    \hline
    #1 & & & & & & & & & \\%
}%

% Environnement titre de fiche de suivi
\define@cmdkey[seq]{titreFiche}{titreSeq}{}%
\define@cmdkey[seq]{titreFiche}{numSeq}{}%
\define@cmdkey[seq]{titreFiche}{classe}{}%
\define@cmdkey[seq]{titreFiche}{etab}{}%
\define@cmdkey[seq]{titreFiche}{annee}{}%
\define@cmdkey[seq]{titreFiche}{enseignant}{}%

\newenvironment{titreFiche}[1][]{%
    \setkeys[seq]{titreFiche}{#1}%
    \ifdocstd
        \pagestyle{fancy}%
        \fancyhf{}%
        \lfoot{Le \today}%
        \rfoot{Page \thepage/\pageref{LastPage}}%
    \fi
    \seqSetColors% Couleurs par défaut (N00/S00 en variables internes → gris)
    \begin{boiteTitreGen}{Séquence \cmdseq@titreFiche@numSeq}%
        \begin{minipage}[c][70pt][c]{\linewidth}%
            \vfill \textit{Fiche de suivi}%
            \ifx\cmdseq@titreFiche@annee\undefined\else
                \hfill \textit{Année \cmdseq@titreFiche@annee}%
            \fi
            \vfill \centering {\Huge \textbf{\cmdseq@titreFiche@titreSeq}}%
            \vfill
            \ifx\cmdseq@titreFiche@enseignant\undefined\else
                \textit{Enseignant : \cmdseq@titreFiche@enseignant}%
            \fi
            \hfill \textit{\cmdseq@titreFiche@etab - \cmdseq@titreFiche@classe}%
        \end{minipage}%
    \end{boiteTitreGen}%
    \par\vspace{10pt}%
    % Migration longtabu → longtblr (tabularray)
    % X[9,l]          : colonne nom élève, 9 parts, alignée à gauche
    % *{9}{X[c]}      : 9 colonnes de maîtrise, centrées
    % vline{2}        : trait vertical avant la 1ère colonne de maîtrise
    % rowsep=5pt      : équivalent \tabulinesep=5pt
    % rowhead=1       : répétition de l'en-tête sur chaque page (= \endhead)
    % hlines          : traits horizontaux automatiques (remplace \tabucline{-})
    \begin{longtblr}[
        caption = {},
    ]{%
        width = \linewidth,
        colspec = { X[9,l] *{9}{X[c]} },
        vline{2},
        rowsep = 5pt,
        rowhead = 1,
        hlines,
        row{1} = {font=\bfseries},
    }%
        & \rotatebox{90}{Niveau}%
        & \rotatebox{90}{Note globale}%
        & \rotatebox{90}{Comportement}%
        & \rotatebox{90}{Restitution cours}%
        & \rotatebox{90}{Restitution exos}%
        & \rotatebox{90}{Éval bilan}%
        & \rotatebox{90}{Série Fondamentale}%
        & \rotatebox{90}{Série Avancée}%
        & \rotatebox{90}{Série Exploration} \\%
}{%
    \end{longtblr}%
}%

%% ------------------------------------------------------------------
%% ---- Décodage et colorisation (maîtrise) -------------------------
%% ------------------------------------------------------------------

% Décodage des séries d'exercices
\newcommand{\seqDecodeSerie}[1]{%
    \ifthenelse{\equal{#1}{AE}}{Révision}{%
        \ifthenelse{\equal{#1}{F}}{Fondamental}{%
            \ifthenelse{\equal{#1}{A}}{Avancé}{%
                \ifthenelse{\equal{#1}{E}}{Exploration}{INCONNU:#1}%
            }%
        }%
    }%
}%

% Mise en couleur des niveaux de maîtrise
\newcommand{\seq@maitrise@colore}[1]{%
    \DTLifeq{#1}{I}{\colorbox{red}{\color{white} \textbf{I}}}{%
        \DTLifeq{#1}{F}{\colorbox{yellow}{\textbf{F}}}{%
            \DTLifeq{#1}{A}{\colorbox{green!60!white}{\textbf{A}}}{%
                \DTLifeq{#1}{E}{\colorbox{green!60!black}{\color{white}\textbf{E}}}{%
                    \DTLifeq{#1}{NE}{NE}{%
                        \DTLifeq{#1}{-}{-}{\colorbox{gray!50!white}{\color{white}\textbf{#1}}}%
                    }%
                }%
            }%
        }%
    }%
}%

% Décodage des niveaux de maîtrise
\newcommand{\seq@maitrise@decode}[1]{%
    \DTLifeq{#1}{I}{\colorbox{red}{\color{white} \textbf{Insuffisant}}}{%
        \DTLifeq{#1}{F}{\colorbox{yellow}{\textbf{À consolider}}}{%
            \DTLifeq{#1}{A}{\colorbox{green!60!white}{\textbf{Satisfaisant}}}{%
                \DTLifeq{#1}{E}{\colorbox{green!60!black}{\color{white}\textbf{Très bon}}}{%
                    \DTLifeq{#1}{NE}{Non évalué}{%
                        \DTLifnumerical{#1}{#1}{%
                            \DTLifeq{#1}{-}{-}{INCONNU:#1}%
                        }%
                    }%
                }%
            }%
        }%
    }%
}%

%% ------------------------------------------------------------------
%% ---- Bilan périodique par élève ----------------------------------
%% ------------------------------------------------------------------

\define@cmdkey[seq]{bilanPeriodiqueParEleve}{periode}{}%
\define@cmdkey[seq]{bilanPeriodiqueParEleve}{annee}{}%
\define@cmdkey[seq]{bilanPeriodiqueParEleve}{niveau}{}%
\define@cmdkey[seq]{bilanPeriodiqueParEleve}{nomClasse}{}%

\newcommand{\seqBilanPeriodiqueParEleve}[1][]{%
    \setkeys[seq]{bilanPeriodiqueParEleve}{#1}%
    \DTLloaddb{sequencesPeriode\cmdseq@bilanPeriodiqueParEleve@periode}{../SequencesDB/\cmdseq@bilanPeriodiqueParEleve@periode -sequences.csv}%
    \DTLloaddb[keys={Numero}]{bilanPeriode\cmdseq@bilanPeriodiqueParEleve@periode}{./\cmdseq@bilanPeriodiqueParEleve@periode Bilan.csv}%
    \gdef\doamp{\gdef\doamp{&}}%
    \fancyhead[C]{\large{Bilan de période~-~\cmdseq@bilanPeriodiqueParEleve@periode}}%
    \fancyhead[R]{\cmdseq@bilanPeriodiqueParEleve@nomClasse~-~\cmdseq@bilanPeriodiqueParEleve@annee}%
    \DTLforeach*{eleves}{\numEleve=Numero,\nomEleve=Nom,\prenomEleve=Prenom}{%
        \begin{minipage}{\linewidth}%
            \hrule%
            \colorbox{blue!30!white}{\textbf{\nomEleve~\prenomEleve}} \medskip \par%
            \DTLforeach*{sequencesPeriode\cmdseq@bilanPeriodiqueParEleve@periode}{\tmp@codesequence=Code,\nomSeq=Nom}{%
                \begin{minipage}{0.6\linewidth} \textbf{\tmp@codesequence}~-~\nomSeq \end{minipage} \hfill%
                \seqDTLfetchsave{\niveauMaitrise}{bilanPeriode\cmdseq@bilanPeriodiqueParEleve@periode}{Numero}{\numEleve}{\tmp@codesequence}%
                \seq@maitrise@decode{\niveauMaitrise} \medskip \par%
            }%
        \end{minipage}%
        \bigskip \par%
    }%
}%

% Les macros seqConfigPlan, seqSetSequence sont supprimées.
% Les alias exercice et serieExos sont supprimés.
% Utiliser : \seqSetCodeNiveau, \seqSetCodeSequence, \seqSetColors,
%            \begin{seqExercice}, \begin{seqSerieExos}.

%% Compatibilité dans le fichier C04_attendus_fin_C03.csv
\newenvironment{prerequisN2}{\begin{seq@prerequis@N2}}{\end{seq@prerequis@N2}}%

%% ------------------------------------------------------------------
%% ---- Bilan par compétences (modèle 2021) -------------------------
%% ------------------------------------------------------------------
%% Macros et environnements du modèle "Bilan des compétences" tel
%% qu'il était utilisé en 2021 (Bilan_T1, différenciation par élève
%% via \DTLforeach sur CSV). Conservés ici pour archivage et pour
%% permettre la recompilation occasionnelle d'anciens documents.
%%
%% Déplacé depuis seqenseigne-core.dtx en v0.13.5.2.
%% Le modèle de référence pour les nouvelles évaluations est défini
%% dans seqenseigne-core-eval.dtx (sans différenciation par élève).
%%
%% Le compteur seqBilanExoNum est désormais déclaré globalement dans
%% seqenseigne-core (v0.13.5.2 étape 2 : centralisation des compteurs).

%% ---- seqBilanItemAlea --------------------------------------------
%    \begin{macro}{\seqBilanItemAlea}%
\newcommand{\seqBilanItemAlea}[3] {%
	\midrule \QcmItemNum
	\FPrandom\hasard \FPupn\varA{3 \hasard{} * 0 trunc}		% entre 0 et 2
	\FPrandom\hasard \FPupn\varB{2 \hasard{} * 0 trunc}		% entre 0 et 1
	\ifthenelse{\varA=0}{%
		\ifthenelse{\varB=0}{#1 & #2 & #3}{#1 & #3 & #2}%
	}{%
		\ifthenelse{\varA=1}{%
			\ifthenelse{\varB=0}{#2 & #1 & #3}{#3 & #1 & #2}%
		}{%
			\ifthenelse{\varB=0}{#2 & #3 & #1}{#3 & #2 & #1}%
		}%
	}%
	\\
}%
%    \end{macro}%
% -------------------------------------------------------------

%% ---- Pour chaque corrigé -----------------------------------------
%    \begin{macro}{\seqBilanCorrigeExo}%
%% Ecrit le corrigé de l'exercice dans un fichier séparé
\newcommand{\seqBilanCorrigeExo}[1]{%
	\scantokens{%
		\begin{tcbverbatimwrite}{Corriges/c\theseqBilanExoNum.tex}%
			#1
		\end{tcbverbatimwrite}%
	}%
}%
%    \end{macro}%
% -------------------------------------------------------------

%% ---- Regroupe les corrigés en fin de livret ----------------------
%    \begin{macro}{\seqAfficheCorrigesBilan}%
%% qui commence sur une page impaire
\newcommand{\seqAfficheCorrigesBilan}{%
	\immediate\closeout\corrfileeval
	\clearpage
	\ifodd\c@page\else
		\null\newpage
		\if@twocolumn\null\newpage\fi
	\fi
	 \clearpage \seqTitreSectionCorrige { Corrigés des exercices de la série exploration } \label {corrigeExos3} 
 \par \nopagebreak \vspace {5pt}\nopagebreak \begin {boiteTitreExo} \par Exercice\;2\par \par \textit {Des caractères sacrés}\end {boiteTitreExo} \par \nopagebreak \vspace {5pt}\nopagebreak   \begin {seqColEnum}[nbCols=1] \item Je sais que $1~302 = 1\times 1~000 + 3 \times 100 + 2 \times 1$. Or, dans le premier nombre en hiéroglyphes, je vois : \begin {seqColItem}[nbCols=1] \item une seule fois le premier hiéroglyphe, \includegraphics [scale=0.2]{N09S01E02_corrige4.png}, donc il vaut 1~000. \item trois fois le deuxième hiéroglyphe, \includegraphics [scale=0.2]{N09S01E02_corrige3.png}, donc il vaut 100. \item deux fois le troisième hiéroglyphe, \includegraphics [scale=0.2]{N09S01E02_corrige1.png}, donc il vaut 1. \end {seqColItem} Je sais que $10~047 = 1 \times 10~000 + 4 \times 10 + 7 \times 1$. Or, dans le deuxième nombre en hiéroglyphes, je vois : \begin {seqColItem}[nbCols=1] \item une seule fois le premier hiéroglyphe, \includegraphics [scale=0.2]{N09S01E02_corrige5.png}, donc il vaut 10~000. \item quatre fois le deuxième hiéroglyphe, \includegraphics [scale=0.2]{N09S01E02_corrige2.png}, donc il vaut 10. \item sept fois le le troisième hiéroglyphe, \includegraphics [scale=0.2]{N09S01E02_corrige1.png} que je connais déjà. \end {seqColItem} Je sais que $2~311~021 = 2 \times 1~000~000 + 3 \times 100~000 + 1 \times 10~000 + 1 \times 1~000 + 2 \times 10 + 1 \times 1$. \newline Or, dans le troisième nombre en hiéroglyphes, je vois : \begin {seqColItem}[nbCols=1] \item deux fois le premier hiéroglyphe, \includegraphics [scale=0.2]{N09S01E02_corrige7.png}, donc il vaut 1~000~000. \item trois fois le deuxième hiéroglyphe, \includegraphics [scale=0.2]{N09S01E02_corrige6.png}, donc il vaut 100~000. \item une seule fois le troisième hiéroglyphe, \includegraphics [scale=0.2]{N09S01E02_corrige5.png} que je connais déjà. \item une seule fois le quatrième hiéroglyphe, \includegraphics [scale=0.2]{N09S01E02_corrige4.png} que je connais déjà. \item deux fois le cinquième hiéroglyphe, \includegraphics [scale=0.2]{N09S01E02_corrige2.png} que je connais déjà. \item une seule fois le sixième hiéroglyphe, \includegraphics [scale=0.2]{N09S01E02_corrige1.png} que je connais déjà. \end {seqColItem} Voici donc dans l'ordre les valeurs des sept hiéroglyphes : \begin {seqColItem}[nbCols=3] \item \includegraphics [scale=0.2]{N09S01E02_corrige1.png} vaut 1. \newline \item \includegraphics [scale=0.2]{N09S01E02_corrige2.png} vaut 10. \newline \item \includegraphics [scale=0.2]{N09S01E02_corrige3.png} vaut 100. \newline \item \includegraphics [scale=0.2]{N09S01E02_corrige4.png} vaut 1~000. \newline \item \includegraphics [scale=0.2]{N09S01E02_corrige5.png} vaut 10~000. \newline \item \includegraphics [scale=0.2]{N09S01E02_corrige6.png} vaut 100~000. \newline \item \includegraphics [scale=0.2]{N09S01E02_corrige7.png} vaut 1~000~000. \end {seqColItem} \item En m'aidant de la question précédente, je peux calculer les nombres pour chaque hiéroglyphe : \newline \includegraphics [scale=0.2]{N09S01E02_enonce4.png} \quad vaut $3 \times 10 + 2 \times 1 = 32$ \newline \includegraphics [scale=0.2]{N09S01E02_enonce5.png} \quad vaut $1 \times 100~000 + 3 \times 10~000 + 2 \times 1 = 130~002$ \newline \includegraphics [scale=0.2]{N09S01E02_enonce6.png} \quad vaut $1 \times 1~000~000 + 2 \times 1~000 + 1 \times 10 = 1~002~010$ \item Je peux écrire $426 = 4 \times 100 + 2 \times 10 + 6 \times 1$ et donc \includegraphics [scale=0.9]{N09S01E02_corrige8.png}. \newline Je peux écrire $527 = 5 \times 100 + 2 \times 10 + 7 \times 1$ et donc \includegraphics [scale=0.9]{N09S01E02_corrige9.png}. \newline Je peux écrire $12~315 = 1 \times 10~000 + 2 \times 1~000 + 3 \times 100 + 1 \times 10 + 5 \times 1$ et donc \includegraphics [scale=0.9]{N09S01E02_corrigeA.png}. \newline Je peux écrire $1~234~000 = 1 \times 1~000~000 + 2 \times 100~000 + 3 \times 10~000 + 4 \times 1~000$ et donc \includegraphics [scale=0.9]{N09S01E02_corrigeB.png}. \newline \end {seqColEnum}  
 
%
	\clearpage
	\ifodd\c@page\else
		\null\newpage
		\if@twocolumn\null\newpage\fi
	\fi

}%
%    \end{macro}%
% -------------------------------------------------------------


%    \begin{macro}{seqBilan}%
%% ---- Options
\define@cmdkey[seq]{bilan}{periode}{}%
\define@cmdkey[seq]{bilan}{annee}{}%
\define@cmdkey[seq]{bilan}{etab}{}%
\define@cmdkey[seq]{bilan}{classe}{}%
\define@cmdkey[seq]{bilan}{eleve}{}%
\newcommand{\seqBilan}[1][]{%
	%% ---- Options
	\setkeys[seq]{bilan}{#1}%
	%% ---- Début du corrigé : titre
	\immediate\openout\corrfileeval=CorrList.tex
	\immediate\write\corrfileeval{%
		\noexpand\cleardoublepage
		\noexpand\begin{boiteTitreGen}%
			\@charlb \eue{\cmdseq@bilan@annee} - \eue{\cmdseq@bilan@periode} \@charrb
			\unexpanded{\begin{minipage}[c][70pt][c]{\linewidth}}%
				\noexpand\vfill \noexpand\centering
				\@charlb \unexpanded{ \Huge \textbf{Bilan des compétences - CORRIG\'E}} \@charrb
				\noexpand\vfill
				\noexpand\textit{\eue{\cmdseq@bilan@eleve}}%
				\noexpand\hfill
				\noexpand\textit{\eue{\cmdseq@bilan@etab} - \eue{\cmdseq@bilan@classe}}%
				\unexpanded{ \end{minipage}	\end{boiteTitreGen} }%
}%
%% ---- Entêtes et pieds de page
\ifdocstd
	\pagestyle{fancy}\fancyhf{}%
	\lfoot{Le \today}\rfoot{Page \thepage/\pageref{LastPage}}%
\fi
%% ---- Couleurs
\seq@setColors{gris}%
%% ---- Boîte de titre
\begin{boiteTitreGen}{\cmdseq@bilan@annee~-~\cmdseq@bilan@periode}%
	\begin{minipage}[c][70pt][c]{\linewidth}%
		\vfill \centering {\Huge \textbf{Bilan des compétences}}%
		\vfill \textit{\cmdseq@bilan@eleve} \hfill
		\textit{\cmdseq@bilan@etab~-~\cmdseq@bilan@classe}%
	\end{minipage}%
\end{boiteTitreGen}%
}%
%    \end{macro}%
% -------------------------------------------------------------

%% ---- Item de liste -----------------------------------------
%    \begin{macro}{\seqBilanObjectifsItem}%
\newcommand{\seqBilanObjectifsItem}[2]{%
	\midrule #1 & #2 \\
}%
%    \end{macro}%

%% ---- Remise à zéro du numéro d'exercice --------------------
%    \begin{macro}{\seqBilanRAZExoNum}%
\newcommand{\seqBilanRAZExoNum}{%
	\setcounter{seqBilanExoNum}{0}%
}%
%    \end{macro}%
% -------------------------------------------------------------

%% ---- Liste des objectifs et niveaux ------------------------------
%    \begin{environment}{seqBilanObjectifs}%
\NewEnviron{seqBilanObjectifs}{%
	\begin{center} \begin{boitePaleVariable}{\linewidth}%
			\begin{tblr}[expand=\BODY]{X[5]X[1.5,c]X[c]}%
				\toprule \large \textbf{Séquence} & \large \textbf{Niveau} & \large \textbf{Validé ?} \\
				\BODY
				\bottomrule
			\end{tblr} \smallskip \newline
			\large \textbf{Consignes}: \newline
			Toutes les réponses doivent être justifiées, 	sauf si une indication contraire est donnée. \newline
			Pour chaque question, si le travail n'est pas terminé, laisser tout de même une trace de la recherche: elle sera prise en compte dans l'évaluation. \smallskip \newline
			\large \textbf{Remarques du correcteur}: \bigskip \bigskip \bigskip
		\end{boitePaleVariable} \end{center}%
}%
%    \end{environment}%
% -------------------------------------------------------------

%% ---- QCM ---------------------------------------------------------
%    \begin{environment}{seqBilanQCM}%
%% ---- Options
\define@cmdkey[seq]{bilanQCM}{nbReps}[3]{}%
\define@cmdkey[seq]{bilanQCM}{largeurQuestion}[3]{}%
\define@cmdkey[seq]{bilanQCM}{largeurRep}[1]{}%
\define@cmdkey[seq]{bilanQCM}{titre}[QCM]{}%
\newenvironment{seqBilanQCM}[1][]{%
	%% ---- Options
	\setkeys[seq]{bilanQCM}{nbReps,largeurQuestion,largeurRep,titre,#1}%
	%% ---- Numéro de l'exercice
	\stepcounter{seqBilanExoNum}%
	%% Note v0.13.5.2 étape 2 : un \setcounter{seqQcmItemNum}{0} apparaissait
	%% ici à l'origine, lié à la macro \QcmItemNum (référencée dans
	%% \seqBilanItemAlea mais jamais définie dans le paquet — vestige
	%% d'une version antérieure). Comme tout le code Bilan 2021 est en
	%% archivage et que les .tex 2019/2021 ne compilent plus avec le
	%% paquet courant, ce setcounter a été retiré sans remplacement.
	%% ---- Début du corrigé du QCM
	\immediate\write\corrfileeval{%
		\unexpanded{\par\nopagebreak\vspace{5pt}\nopagebreak}%
		\unexpanded{\begin{boiteTitreSection}}%
				\noexpand\begin{minipage}[c][20pt][c]{\noexpand\linewidth} \noexpand\Large
					\noexpand\textbf{ Exercice \Roman{seqBilanExoNum} : \eue{\cmdseq@bilanQCM@titre} }%
					\unexpanded{\end{minipage} \end{boiteTitreSection}}%
		\unexpanded{\par\nopagebreak\vspace{5pt}\nopagebreak}%
		\noexpand\input{Corriges/c\theseqBilanExoNum.tex}%
	}%
	%% ---- Boite de titre
	\par\vspace{10pt}%
	\begin{boiteTitreSection}%
		\begin{minipage}[c][20pt][c]{\linewidth}%
			\Large \textbf{Exercice \Roman{seqBilanExoNum} : \cmdseq@bilanQCM@titre}%
		\end{minipage}%
	\end{boiteTitreSection}%
	\par\vspace{5pt}%
	Cet exercice est un questionnaire à choix multiples :
	pour chaque question, il y a une ou plusieurs réponse(s) correcte(s).
	Pour répondre, indiquez le numéro de la question et les codes (A, B, C)
	correspondant aux affirmation(s) juste(s) sur votre copie. \newline
	\centerline{\textbf{Aucune justification n'est attendue pour ce QCM.}}%
	\setcounter{seq@qcm@alphcode}{40}%
	\renewcommand\seq@qcm@repCodes{}%
	\forloop{seq@qcm@iter}{0}{\value{seq@qcm@iter} < \cmdseq@bilanQCM@nbReps}%
	{\stepcounter{seq@qcm@alphcode} \eappto\seq@qcm@repCodes{ & \char"\theseq@qcm@alphcode}}%
	\xltabular{0.95\linewidth}{L{0.45\linewidth} *{3}{C{0.15\linewidth}}}%
	\toprule
	\centering\textbf{Questions} & \multicolumn{\cmdseq@bilanQCM@nbReps}{c}{\textbf{Affirmations}} \\
	\seq@qcm@repCodes \\
	\endhead
}{%
	\endxltabular
}%
%    \end{environment}%
% -------------------------------------------------------------

%% ---- Exercice autre que QCM --------------------------------------
%    \begin{environment}{seqBilanExo}%
%% ---- Options
\define@cmdkey[seq]{seqBilanExo}{titre}{}%
\newenvironment{seqBilanExo}[1][]{%
	%% ---- Options
	\setkeys[seq]{seqBilanExo}{#1}%
	%% ---- Numéro de l'exercice
	\stepcounter{seqBilanExoNum}%
	%% ---- Début du corrigé
	\immediate\write\corrfileeval{%
		\unexpanded{\par\nopagebreak\vspace{5pt}\nopagebreak}%
		\unexpanded{\begin{boiteTitreSection}}%
				\noexpand\begin{minipage}[c][20pt][c]{\noexpand\linewidth}%
					\noexpand\Large \noexpand\textbf{ Exercice \Roman{seqBilanExoNum}%
						\ifx\cmdseq@seqBilanExo@titre\undefined\else : \eue{\cmdseq@seqBilanExo@titre} \fi }%
					\unexpanded{\end{minipage} \end{boiteTitreSection}}%
		\unexpanded{\par\nopagebreak\vspace{5pt}\nopagebreak}%
		\noexpand\input{Corriges/c\theseqBilanExoNum.tex}%
	}%
	%% ---- Boite de titre
	\par\vspace{10pt}%
	\begin{boiteTitreSection}%
		\begin{minipage}[c][20pt][c]{\linewidth}%
			\Large \textbf{Exercice \Roman{seqBilanExoNum}%
				\ifx\cmdseq@seqBilanExo@titre\undefined\else : \cmdseq@seqBilanExo@titre\fi }%
		\end{minipage}%
	\end{boiteTitreSection}%
	\par\vspace{5pt}%
}{}%
%    \end{environment}%
% -------------------------------------------------------------

% Fin du module legacy
%</package>
%
% \endinput